On MW-32 the result is unambiguous. Trained on exact samples, iDEM learns a score whose relative error is $25.0\pm2.6$ at $\sigma_t=1$ and $35.0\pm2.3$ at $\sigma_t=2$ (3 seeds), against $0.140\pm0.003$ and $0.074\pm0.004$ for AEM. Ten times more proposals barely help iDEM: with $K=1000$ (one seed) its error is still $9.8$--$27.6$. This is what one expects from an SNIS bias that decays like $1/K$ with a constant that is exponential in $d$. With a perfect buffer, AEM also gets the mode weights right: in seed 0 the fraction of double-well coordinates in the major well is $0.837$, against a true $0.844$ (iDEM: $0.898$). So the mode-weight deficits of \cref{tab:main} come from the self-generated buffer, as \cref{thm:contraction} predicts, and not from the anchored target.
On GMM-10, where the TSI/DSI combination is exact within each unit-variance component, the two methods learn similar scores at $\sigma_t\ge0.3$ (differences within one standard deviation); AEM is more accurate at $\sigma_t\le0.1$ in the mean.
On the molecular targets we rank the models by held-out denoising loss, which by \cref{prop:diag}(ii) differs between two models by exactly the difference of their score errors (normalized as $\tilde L=\sigma_t^2L/d_{\mathrm{eff}}$; the held-out configurations are disjoint from every seed's buffer). On DW-4 the AEM model has the lower score error at every noise level and in every seed at $\sigma_t\ge0.3$, by a margin that grows from $0.02$ to $10.9$ (\cref{tab:diag}); its samples also have a lower interatomic-distance $\mathcal{W}_1$ in every seed ($0.035$--$0.101$ vs.\ $0.085$--$0.250$). On LJ-13, AEM is better at $\sigma_t\in\{1,2\}$ in all seeds, but slightly worse in the mean at $\sigma_t\le0.3$ (not consistent across seeds), and the gaps are small compared with DW-4. The oracle-trained LJ-13 models of both methods produce some very high-energy samples, so we do not use their sample metrics.
